\documentclass[10pt,a4paper]{amsart}
\usepackage[margin=19mm]{geometry}
\usepackage{amsmath,amssymb,mathtools}
\usepackage[scr=rsfs,cal=euler]{mathalfa}
\usepackage{xfrac}
\usepackage[unicode,hidelinks,bookmarks=false]{hyperref}
\newcommand{\ie}{i.e.\ }
\newcommand{\cf}{cf.\ }
\newcommand{\resp}{respectively\ }
\newcommand{\Subsection}{\subsection}
\providecommand{\nobreakdash}{\nobreak\hbox{-}}
\setlength{\emergencystretch}{2em}
\multlinegap0pt
\usepackage{amssymb}
\usepackage{cancel}
\usepackage[shortlabels]{enumitem}
\usepackage{mathtools}
\usepackage{xcolor}
\usepackage{enumerate}
\usepackage{float}
\newtheorem{definition}{Definition}
\newtheorem{theorem}{Theorem}
\newcommand{\beq}{\begin{equation}}
\newcommand{\bs}{\boldsymbol}
\newcommand{\eeq}{\end{equation}}
\newcommand{\rd}{\mathrm{d}}
\title{Jacobi-Haantjes manifolds, Integrability \\ and Dissipative Mechanical Systems}
\author{Rafael Azuaje, Piergiulio Tempesta}
\date{Source: arXiv:2507.11715v3 (CC BY 4.0)}
\begin{document}
\maketitle

\noindent\emph{Source and licence.} Jacobi-Haantjes manifolds, Integrability \\ and Dissipative Mechanical Systems. Version arXiv:2507.11715v3 (CC BY 4.0), distributed by arXiv under the Creative Commons Attribution 4.0 International licence, http://creativecommons.org/licenses/by/4.0/, as recorded in the arXiv licence metadata for this exact version and shown on the abstract page https://arxiv.org/abs/2507.11715v3. Full licence notice: \texttt{licenses/arxiv-2507.11715-CC-BY-4.0.txt}.\par\smallskip

\noindent\textit{Editorial scope.} Counted unit: the article's theorem themain (source0.tex lines 607--611). The article's complete original proof of it is delivered with it (source0.tex lines 613--719, one \texttt{proof} environment of 107 lines). No mathematics is written, repaired or re-derived: the statement and the proof are the article's own lines, in the article's own notation, and no retained line is substituted or re-flowed.  Retained uncounted as the source-local context the proof needs: lines 586--593; lines 598--604.  Nothing was omitted from any retained span, so the package carries no omission record.  No retained line contains a citation, so the wrapper carries no bibliography and no bibliography entry.  No retained line contains a figure environment, an image-inclusion command or drawing code, so no image is delivered and the package declares no asset dependency.  Editorial formatting only, in addition: an AMS \texttt{amsart} preamble that restates the article's own declarations and macros used by the retained lines, namely the environments \texttt{definition}, \texttt{theorem} on separate counters, together with the standard packages those lines need.  The retained environments print in this document as Definition 1, Theorem 1 in order of appearance, whereas the article prints the same items with the numbering of its own full document.  The complete native source of the article as distributed in the arXiv e-print for this exact version is bundled unchanged as \texttt{sources/181649/source0.tex}.
\begin{equation}
\label{systemcc}
\left\lbrace \begin{array}{c}
\Lambda(\bs{K}_{i}^{T}\alpha,\beta)+\alpha(Y_{i})\beta(E)=\Lambda(\alpha,\bs{K}_{i}^{T}\beta)-\beta(Y_{i})\alpha(E), \\ 
\Lambda(\gamma_{i},\alpha)=(\bs{K}_{i}^{T}\alpha)(E)-k_{i}\alpha(E), \\ 
\gamma_{i}(E)=0,
\end{array} \right.
\end{equation}

\begin{definition}
Let $\mathscr{H}$ be an extended Haantjes algebra on $M$ of rank $m$. We shall say that a function $H\in C^{\infty}(M)$ generates an extended Haantjes chain $\mathscr{C}$ on $M$ of length $m$, if  there exists a distinguished basis $\{\bs{\mathcal{K}}_1,\ldots, \bs{\mathcal{K}}_m\}$ of $\mathscr{H}$ and a set of functions $H_{1},\ldots,H_{m}\in C^{\infty}(M)$ such that
 \beq \label{eqHCportuguesa}
 (\rd H_{i}, H_{i})= \bs{\mathcal{K}}_i^{T}(\rd H,H), \qquad i=1,\ldots, m.
 \eeq
The functions $H_i$ are said to be the potential functions of the chain.
\end{definition}

\begin{theorem}
\label{themain}
Let $(M,\Lambda,E,\mathscr{H})$ be a Jacobi-Haantjes manifold of class $m$. If $H\in C^{\infty}(M)$ generates an extended Haantjes chain $\mathscr{C}$ on $M$ of length $m$,
then the potential functions $H_{1},\ldots,H_{m}$ of the chain are independent functions pairwise in involution such that $\lbrace H_{i},H\rbrace=0$. 
\end{theorem}

\begin{proof}
Let $H_{1},\ldots,H_{m}$ be the potential functions of the chain $\mathscr{C}$ and $\lbrace \bs{\mathcal{K}}_{1},\ldots,\bs{\mathcal{K}}_{m}\rbrace$ be a distinguished basis of $\mathscr{H}$ such that equation (\ref{eqHCportuguesa}) is satisfied. 
As before, $\bs{\mathcal{K}}_{i}=(\bs{K}_{i},Y_{i},\gamma_{i},k_{i})$, where $\bs{K}_{i}$ is a (1,1)-tensor field on $M$, $Y_{i}\in \mathfrak{X}(M)$, $\gamma_{i}\in \Omega^{1}(M)$ and $k_{i}\in C^{\infty}(M)$. Thus, for all $(X,f)\in\mathfrak{X}(M)\times C^{\infty}(M)$ we obtain
\begin{equation}
\begin{split}
\left( \bs{\mathcal{K}}_{i}^{T}(\rd H,H)\right)(X,f)&=(\rd H,H)\left( \bs{\mathcal{K}}_{i}(X,f)\right)\\
&=(\rd H,H)(\bs{K}_{i}X+fY_{i},\gamma_{i}X+k_{i}f)\\
&=\rd H(\bs{K}_{i}X+fY_{i})+(\gamma_{i}X+k_{i}f)H\\
&=(\bs{K}_{i}^{T}\rd H+H\gamma_{i})X+(Y_{i}H+Hk_{i})f.
\end{split}
\end{equation}
Besides,
\begin{equation}
(\rd H_{i},H_{i})(X,f)=\rd H_{i}(X)+H_{i}f,
\end{equation}
therefore
\begin{equation}
\label{dHi,Hi}
\rd H_{i}=\bs{K}_{i}^{T}\rd H+H\gamma_{i}\quad\textit{and}\quad H_{i}=Y_{i}H+Hk_{i}.
\end{equation}

First, we consider the commutator 
\begin{equation}
\begin{split}
\lbrace H_{i},H\rbrace &=\Lambda(\rd H_{i},\rd H)+H_{i}EH-HEH_{i}\\
&=\Lambda(\bs{K}_{i}^{T}\rd H+H\gamma_{i},\rd H)+H_{i}EH-HEH_{i}\\
&=\Lambda(\bs{K}_{i}^{T}\rd H,\rd H)+H\Lambda(\gamma_{i},\rd H)+H_{i}EH-HEH_{i}\\
&=\Lambda(\bs{K}_{i}^{T}\rd H,\rd H)+H((\bs{K}_{i}^{T}\rd H)(E)-k_{i}(\rd H)(E))+H_{i}EH-HEH_{i}\\
&=\Lambda(\bs{K}_{i}^{T}\rd H,\rd H)+H(\bs{K}_{i}^{T}\rd H)(E)-Hk_{i}EH+H_{i}EH-HEH_{i}.
\end{split}
\end{equation}
Taking into account that
\begin{equation}
(\bs{K}_{i}^{T}\rd H)(E)=(\rd H_{i}-H\gamma_{i})(E)=\rd H_{i}-H\gamma_{i}(E)=EH_{i};
\end{equation}
we obtain
\begin{equation}
\begin{split}
\lbrace H_{i},H\rbrace &=\Lambda(\bs{K}_{i}^{T}\rd H,\rd H)+HEH_{i}-Hk_{i}EH+H_{i}EH-HEH_{i}\\
&=\Lambda(\bs{K}_{i}^{T}\rd H,\rd H)-Hk_{i}EH+H_{i}EH\\
&=\Lambda(\bs{K}_{i}^{T}\rd H,\rd H)-(H_{i}-Y_{i}H)EH+H_{i}EH\\
&=\Lambda(\bs{K}_{i}^{T}\rd H,\rd H)+Y_{i}HEH.
\end{split}
\end{equation}
From the first equation of system (\ref{systemcc}) we deduce
\begin{equation}
\begin{split}
\Lambda(\bs{K}_{i}^{T}\rd H,\rd H)+Y_{i}HEH&=\Lambda(\rd H,\bs{K}_{i}^{T}\rd H)-Y_{i}HEH\\
&=-\Lambda(\bs{K}_{i}^{T}\rd H,\rd H)-Y_{i}HEH,
\end{split}
\end{equation}
which implies that $$\Lambda(\bs{K}_{i}^{T}\rd H,\rd H)+Y_{i}HEH=0.$$  Consequently,
\begin{equation}
\lbrace H_{i},H\rbrace=0.
\end{equation}
Let us prove now that $\lbrace H_{i},H_{j}\rbrace=0$. To this aim, observe that
\begin{equation}
\begin{split}
\lbrace H_{i},H_{j}\rbrace &=\Lambda(\rd H_{i},\rd H_{j})+H_{i}EH_{j}-H_{j}EH_{i}\\
&=\Lambda(\bs{K}_{i}^{T}\rd H+H\gamma_{i},\bs{K}_{j}^{T}\rd H+H\gamma_{j})+H_{i}EH_{j}-H_{j}EH_{i}\\
&=\Lambda(\bs{K}_{i}^{T}\rd H,\bs{K}_{j}^{T}\rd H)+H\Lambda(\bs{K}_{i}^{T}\rd H,\gamma_{j})+H\Lambda(\gamma_{i},\bs{K}_{j}^{T}\rd H)\\
&\quad+H^{2}\Lambda(\gamma_{i},\gamma_{j})+H_{i}EH_{j}-H_{j}EH_{i}\\
&=\Lambda(\bs{K}_{i}^{T}\rd H,\bs{K}_{j}^{T}\rd H)+H(k_{j}(\bs{K}_{i}^{T}\rd H)(E)-(\bs{K}_{j}^{T}\bs{K}_{i}^{T}\rd H)(E))\\
&\quad+H((\bs{K}_{i}^{T}\bs{K}_{j}^{T}\rd H)(E)-k_{i}(\bs{K}_{j}^{T}\rd H)(E))+H^{2}((\bs{K}_{i}^{T}\gamma_{j})(E)\\
&\quad-k_{i}\gamma_{j}(E))+H_{i}EH_{j}-H_{j}EH_{i}\\
&=\Lambda(\bs{K}_{i}^{T}\rd H,\bs{K}_{j}^{T}\rd H)+H(k_{j}EH_{i}-k_{i}EH_{j}+\rd H(\bs{K}_{j}\bs{K}_{i}E-\bs{K}_{i}\bs{K}_{j}E))\\
&\quad+H^{2}\gamma_{j}(\bs{K}_{i}E)+H_{i}EH_{j}-H_{j}EH_{i}\\
&=\Lambda(\bs{K}_{i}^{T}\rd H,\bs{K}_{j}^{T}\rd H)+H_{j}EH_{i}-Y_{j}HRH_{i}-H_{i}EH_{j}+Y_{i}HEH_{j}\\
&\quad+H\rd H(\bs{K}_{j}\bs{K}_{i}E-\bs{K}_{i}\bs{K}_{j}E))+H^{2}\gamma_{j}(\bs{K}_{i}E)+H_{i}EH_{j}-H_{j}EH_{i}\\
&=\Lambda(\bs{K}_{i}^{T}\rd H,\bs{K}_{j}^{T}\rd H)-Y_{j}HRH_{i}+Y_{i}HEH_{j}\\
&\quad+H\rd H(\bs{K}_{j}\bs{K}_{i}E-\bs{K}_{i}\bs{K}_{j}E))+H^{2}\gamma_{j}(\bs{K}_{i}E).
\end{split}
\end{equation}
Since $\mathscr{H}$ is an Abelian algebra, we necessarily have $\bs{\mathcal{K}}_{i}\bs{\mathcal{K}}_{j}=\bs{\mathcal{K}}_{j}\bs{\mathcal{K}}_{i}$. For all $(X,f)\in\mathfrak{X}(M)\times C^{\infty}(M)$ we obtain
\begin{equation}
\bs{\mathcal{K}}_{i}\bs{\mathcal{K}}_{j}(X,f)=(\bs{K}_{i}\bs{K}_{j}X+Y_{i}\gamma_{j}X+f\bs{K}_{i}Y_{j}+fk_{j}Y_{i},(\bs{K}_{j}^{T}\gamma_{i})(X)+k_{i}\gamma_{j}X+f\gamma_{i}Y_{j}+fk_{i}k_{j}).
\end{equation}
Explicitly, 
\begin{equation}
\bs{\mathcal{K}}_{i}\bs{\mathcal{K}}_{j}=(\bs{K}_{i}\bs{K}_{j}+Y_{i}\otimes\gamma_{j},\bs{K}_{i}Y_{j}+k_{j}Y_{i},\bs{K}_{j}^{T}\gamma_{i}+k_{i}\gamma_{j},\gamma_{i}Y_{j}+k_{i}k_{j}),
\end{equation}
where $Y_{i}\otimes\gamma_{j}$ is the (1,1)-tensor field on $M$ defined by $(Y_{i}\otimes\gamma_{j})(X):=Y_{i}\gamma_{j}X$. By taking $X=E$ we obtain
\begin{equation} \label{eq:4.16}
\left\lbrace \begin{array}{c}
\bs{K}_{i}\bs{K}_{j}E=\bs{K}_{j}\bs{K}_{i}E,\\
\gamma_{i}(\bs{K}_{j}E)=\gamma_{j}(\bs{K}_{i}E),\\
\gamma_{i}Y_{j}=\gamma_{j}Y_{i}.
\end{array} \right. 
\end{equation}
Now, as $\lbrace H_{i},H_{j}\rbrace=-\lbrace H_{j},H_{i}\rbrace$, we have that  $\gamma_{i}(\bs{K}_{j}E)=-\gamma_{j}(\bs{K}_{i}E)$. Combining this relation with the second one of system \eqref{eq:4.16}, we deduce that $\gamma_{i}(\bs{K}_{j}E)=0$. Consequently, the expression of the commutator $\lbrace H_{i},H_{j}\rbrace$ reduces to
\begin{equation}
\lbrace H_{i},H_{j}\rbrace=\Lambda(\bs{K}_{i}^{T}\rd H,\bs{K}_{j}^{T}\rd H)-Y_{j}HRH_{i}+Y_{i}HEH_{j}.
\end{equation}

On the other hand, $\bs{\mathcal{K}}_{i}\bs{\mathcal{K}}_{j}=: \bs{\mathcal{K}}_{s}$ is an element of the algebra $\mathscr{H}$. Thus,
\begin{equation}
\Lambda(\bs{K}_{s}^{T}\rd H,\rd H)+Y_{s}HEH=0.
\end{equation}
Being $\bs{K}_{s}=\bs{K}_{i}\bs{K}_{j}+Y_{i}\otimes\gamma_{j}$ and $Y_{s}=\bs{K}_{i}Y_{j}+k_{j}Y_{i}$, we can easily prove that
\begin{equation}
\Lambda(\bs{K}_{s}^{T}\rd H,\rd H)+Y_{s}HEH=\Lambda(\bs{K}_{i}^{T}\rd H,\bs{K}_{j}^{T}\rd H)-Y_{j}HRH_{i}+Y_{i}HEH_{j},
\end{equation}
i.e.,
\begin{equation}
\lbrace H_{i},H_{j}\rbrace=\Lambda(\bs{K}_{s}^{T}\rd H,\rd H)+Y_{s}HEH=0.
\end{equation}
\end{proof}

\end{document}
